% ==================================================================
%  DIMENSIONALITY REDUCTION TECHNIQUE TEMPLATE --- copy me, do not
%  edit in place.
%
%    cp template/dimreduction_template.tex \
%       documentation/sections/dimensionality/<linear|nonlinear>/NN-<name>.tex
%
%  Then: replace <NAME>, replace the label suffix `template' with the
%  technique's short name, and add one \input line under the matching
%  subsection (6.3 linear, 6.4 nonlinear) in dim_main.tex.
%
%  This file is NOT included in the build.
%
%  Levels: a technique is a \subsubsection sitting under its subsection
%  in dim_main.tex, so the parts below are \paragraph.  Keep their
%  order unchanged in every copy.
%
%  Differs from the clustering method template in three parts that have
%  no counterpart there --- `Out-of-sample and inverse mapping',
%  `Choosing the number of components', and a Results part reporting
%  both structure preservation and the downstream effect on clustering.
%  A reduction is never the deliverable; what it does to the clustering
%  that follows it is.
%
%  Code counterpart: xxcluster/dim_red/<linear|nonlinear>/<name>.py.
%  Parts 2, 4 and 7 must agree with that class's Capabilities
%  declaration --- see CONTRIBUTING.md, Part 2.
% ==================================================================
\subsubsection{<NAME>}
\label{sec:dr:template}

\paragraph{Overview}
\label{sec:dr:template:overview}
\placeholder{What the technique does, in one paragraph: what it
maps from, what it maps to, and what it tries to preserve in doing so.
State whether it is linear or nonlinear, and for a nonlinear technique,
its relation to the manifold hypothesis of
Section~\ref{sec:dimensionality}.}

\paragraph{Assumptions and applicability}
\label{sec:dr:template:assumptions}
\placeholder{What the technique assumes about the data --- linearity,
local versus global structure, density, scale, noise --- and the
conditions under which it is a sensible choice for the data described
in Section~\ref{sec:data}. State explicitly whether it is intended for
visualisation, for reducing input to a clustering method, or for both:
these are different jobs and a technique good at one may distort what
the other needs.}

\paragraph{Formulation}
\label{sec:dr:template:formulation}
\placeholder{The objective being optimised and the mapping obtained,
using the symbols declared in the notation table. Cite the original
source.}

\paragraph{Out-of-sample and inverse mapping}
\label{sec:dr:template:mapping}
\placeholder{Two questions, answered explicitly, because the rest of the
pipeline depends on both.

First: is the technique inductive or transductive? That is, does it
learn a mapping that applies to observations unseen during fitting, or
does it embed only the sample it was given? A transductive technique
cannot precede a clustering step in a pipeline that will later see new
data.

Second: does an inverse mapping exist? If so, a component can be read
back in terms of the measured features; if not, say so, and say where
the interpretation of the reduced space will come from instead.}

\paragraph{Hyperparameters and tuning}
\label{sec:dr:template:hyperparameters}
\placeholder{Each hyperparameter, the range searched, and the value
selected --- with the reason. For neighbourhood-based techniques, state
what the neighbourhood size does to the balance between local and
global structure, since it is the parameter most likely to change the
conclusions drawn from a figure.}

\paragraph{Choosing the number of components}
\label{sec:dr:template:ncomponents}
\placeholder{How the target dimension was chosen, and on what evidence:
an explained-variance criterion, an intrinsic dimension estimate (see
Section~\ref{sec:dimensionality}), a reconstruction error, or the
requirements of a downstream method. Choosing two because two is what
plots is a legitimate choice for a figure and not for an analysis ---
if that is the reason, say so.}

\paragraph{Complexity}
\label{sec:dr:template:complexity}
\placeholder{Time and space complexity, and what that means at
\projectname{} data volumes. Note any step that is quadratic in the
number of observations, since that is what decides whether the
technique is usable on the full dataset or only on a sample.}

\paragraph{Application to \projectname{} data}
\label{sec:dr:template:application}
\placeholder{Implementation: library and version, random seeds, and how
the run is reproduced. Then any deviation from the shared preprocessing
pipeline of Section~\ref{sec:data:preprocessing}, and why it was
necessary --- scaling in particular, since most of these techniques are
not scale-invariant.}

\paragraph{Results}
\label{sec:dr:template:results}
\placeholder{Two things, in this order.

Structure preservation: how faithfully the embedding represents the
input, reported quantitatively --- explained variance, trustworthiness,
stress, or reconstruction error, whichever the technique admits --- and
not by appeal to the appearance of a scatter plot.

Effect on downstream clustering: the validity measures of
Section~\ref{sec:methodology:metrics} for the same clustering method
fitted with and without this reduction. This is the part that decides
whether the technique earns its place in the pipeline.}

\paragraph{Limitations}
\label{sec:dr:template:limitations}
\placeholder{Where this technique distorted the data, or would be
expected to. Distinguish what was observed on our data from what is
known from the literature. For nonlinear techniques, state plainly which
properties of the embedding must not be read as properties of the data
--- typically inter-cluster distances, relative cluster sizes, and
apparent density.}

\paragraph{Summary}
\label{sec:dr:template:summary}
\placeholder{Two or three sentences: the verdict on this technique, and
whether it is recommended for visualisation, for the clustering
pipeline, for both, or for neither. Phrased so it can be carried
directly into Section~\ref{sec:comparison}.}
